% Generated by tools/textbook/export_listings.js from tools/decks/mastery.js and tools/concepts/registry.json: do not edit.
\begin{chapstatus}{Stato di apprendimento di questo capitolo}
\masterychip{now}{DA COMPRENDERE OGGI}\enspace le maschere di bit; i codici; gli array in memoria; gli ingressi con pull-up; l'algebra di Boole; la virgola fissa; i multiplexer; istruzioni con o senza segno; comparatori e parità; gli argomenti in a0 ... a7.\par\smallskip
\masterychip{use}{DA UTILIZZARE OGGI}\enspace la convenzione di chiamata (Capitolo~\ref{ch:processor}) e, usati senza essere studiati nel corso, gli strumenti di PRG32, le chiamate del runtime PRG32, il debugger di PRG32-QT.\par\smallskip
\masterychip{later}{DA COMPRENDERE PIÙ AVANTI}\enspace niente.
\end{chapstatus}
\begin{chapstatus}{Percorso del capitolo}
Alla fine del capitolo saprete:
\begin{enumerate}[nosep,leftmargin=2em]
\item leggere i pulsanti con le maschere e rifiutare un'inversione a U con un solo XOR
\item confrontare posizioni con XOR e OR per riconoscere il cibo
\item tenere e mostrare un punteggio
\end{enumerate}
Si appoggia su: le maschere di bit (Capitolo~\ref{ch:boole1}); i codici (Capitolo~\ref{ch:combinational}); gli array in memoria (Capitolo~\ref{ch:information}); gli ingressi con pull-up (Capitolo~\ref{ch:electronics}); l'algebra di Boole (Capitolo~\ref{ch:boole1}). Prossimo passo: Capitolo~\ref{ch:sequential}.
\par\smallskip\textbf{Guardando indietro, guardando avanti.} Nella Lezione 07 una funzione che ne chiama un'altra salvava ra in un frame di 16 byte. Oggi le funzioni di Verm1 si chiamano più in profondità: ciascuna salva ra una volta, all'ingresso, e lo ripristina prima di ret. La Lezione 18 enuncia le regole: quali registri una chiamata può cambiare e quali deve conservare. \emph{Richiamo:} Una funzione chiama altre due funzioni. Quante volte deve salvare ra, e dove?
\end{chapstatus}
